\section{Conclusion}
\label{sec:conclusion}

This paper contributes an open, auditable firm-level accounting microsimulation for
analysing the UK VAT registration threshold, released in full---including its
synthetic-data generator---for replication. Its value is as transparent policy-scenario
infrastructure rather than a structural identification result: I make no claim to
recover policy-invariant firm primitives, unlike \citet{garicano2016} and
\citet{gourioroys2014}. The analysis is anchored to the 2023--24 data year, in
which the statutory threshold was \pounds85{,}000, with a second 2024--25
population (threshold \pounds90{,}000) for the forward sweep.
Table~\ref{tab:summary} collects the headline objects.

\begin{table}[t]
\centering
\caption{Summary of headline estimates, consolidating results defined in the
body sections cross-referenced below. Rows 1--2 are 2025--26 figures (the
anchor on the 2023--24 population, the sweep on the 2024--25 population at its
\pounds90{,}000 baseline); the remaining rows refer to the \pounds85{,}000
data-year baseline.}
\label{tab:summary}
\small
\begin{tabularx}{\textwidth}{@{}>{\raggedright\arraybackslash}p{0.30\textwidth}>{\raggedright\arraybackslash}p{0.19\textwidth}Xl@{}}
\toprule
Object & Definition & Value & Source \\
\midrule
Anchor, \pounds85k to \pounds90k, 2025--26 & dereg.\ / 43\% retention & $-\pounds234$m / $-\pounds134$m (HMRC $-\pounds185$m) & \S\ref{sec:static} \\
Sweep, \pounds90k baseline & \pounds70k to \pounds120k & $+\pounds0.19$bn to $-\pounds1.19$bn & Table~\ref{tab:static_revenue} \\
Dominated region & $a=T^{*}\tau/(1-\tau)$ & \pounds21{,}250 & \S\ref{ssec:dominated} \\
In-scope firms in dominated band & weighted & $\approx$143{,}000 & \S\ref{ssec:dominated} \\
Raise to \pounds100k & static $=$ behavioural & $-\pounds742$m & \S\ref{sec:static}, \S\ref{sec:behavioural} \\
Taper, \pounds85k--\pounds141.7k & linear / constant $50\%$ & $-\pounds1{,}406$m / $-\pounds968$m & \S\ref{ssec:schedule_costs} \\
Reduced rate, \pounds85k--\pounds105k & $10\%$ / $15\%$ & $-\pounds446$m / $-\pounds223$m & \S\ref{ssec:schedule_costs} \\
Bunching readback, 2023--24 & $E$; $b_{\mathrm{LLAT}}$ & $5{,}498$; $5.357$ (target-inherited) & \S\ref{sec:bunching} \\
\bottomrule
\end{tabularx}
\end{table}

\paragraph{Findings.} On the common \pounds85{,}000 baseline, raising the
threshold to \pounds100{,}000 costs \pounds742m, a taper across the
\pounds85{,}000--\pounds141{,}667 band costs \pounds1{,}406m (linear phase-in) or
\pounds968m (constant 50\% marginal rate), and the 10\% and 15\%
reduced-rate bands cost \pounds446m and \pounds223m. The anchor reform
brackets HMRC's published costing between the full-deregistration and 43\%
retention conventions once releases are limited by the deregistration
threshold; it is an internal consistency check rather than external
validation, and the registration behaviour of released firms remains the
decisive convention.

The notch creates an exact \pounds21{,}250 dominated region. Raising the
threshold relocates and widens it, and a reduced-rate band divides it between two notches;
only the continuous taper removes it. Intensive-margin responses change the
reduced-rate costings only modestly under the assumed elasticities, while a
threshold rise has no intensive-margin revenue offset in this model.

The bunching exercise supplies a methodological warning, not a behavioural
estimate. Its 2023--24 profile is inherited from OBR calibration targets, its
degree--window estimates are unstable and fitted missing mass is zero. The
administrative evidence in \citet{liuetal2021} therefore remains the relevant
evidence on firm behaviour.

\paragraph{Limitations and future work.} The model omits location choice,
voluntary-registration choice, compliance costs, sectoral incidence,
pass-through, and general-equilibrium effects. Its population is the ONS
VAT/PAYE frame with an HMRC-calibrated registered subset: threshold cuts are
therefore costed only on frame enterprises, most of which are already
registered, and omit the unregistered sole traders outside the frame that a
lower threshold would mainly reach; extending the frame to the DBT business
population is the natural next build. A fixed cross-section also has no
registration response or entry margin, so the anchor brackets HMRC's profile
only through explicit deregistration conventions. The nominal-growth factors
used for ageing are assumptions recorded in the replication repository. It also cannot capture persistent
growth effects extending beyond the cutoff, as documented for a different
payroll-tax notch by \citet{jysmabenzartiharju2026}. Those estimates do not
transfer directly to UK VAT. Administrative firm-level validation, an explicit
registration model, and welfare accounting are the main priorities for future
work.
